% !PRINTED 215
\weberpage{215}
%<<CONT

\weq{16}{(\gamma\xi+\delta)^{n-2}\, T\!\left(\frac{\alpha\xi+\beta}{\gamma\xi+\delta}\right) = r^{n-2}\Theta(\xi),}

in which the connection between the coefficients $t$ and $\tau$
is expressed by the symbolic equations (13), and is therefore the
same as in the functions $Y$ and $H$.

\gesperrt{Accordingly, $r^{n-2}\Theta(\xi)$ is the transformation of the
function $(\gamma\xi+\delta)^{n-2}T(z)$ by the same linear substitution
by which $(\gamma\xi+\delta)^n f(z)$ passes over into $\varphi(\xi)$.}

If we denote the coefficients of this transformed function,
that is, the quantities $r^{n-2}\tau_k$, by $t'$, then equation (12),
which is linear with respect to $\tau$, yields

\weq{17}{H(t', \xi) = r^{n-1}Y(t, x).}

Now from $Y(t, x)$ there arise $n$ functions when one puts
for $x$ the $n$ roots of $f(x)$. Every homogeneous, rational,
symmetric function of these $n$ functions is rationally expressible
in terms of the coefficients of $f(x)$; and if we therefore denote such
a function by $K(t, a)$, then equation (17) yields, if we denote
the coefficients of $\varphi(\xi)$ by $a'$,

\weq{18}{K(t', a') = r^{\nu(n-1)}K(t, a),}

where $\nu$ denotes the degree of the symmetric function. With
this the beautiful theorem of \gesperrt{Hermite} is proved:

\gesperrt{The coefficients in the equation $f(x) = 0$ transformed
by the Tschirnhausen transformation}

\[
y = Y(t, x)
\]

\gesperrt{and all symmetric functions of the $n$ values $y$ are
simultaneous invariants of the two functions}

\[
f(z), \quad T(z).
\]

\begin{center}
\S.\,70.\\[2pt]
\gesperrt{Developments concerning the Hermite theorem.}
\end{center}

From the considerations of the last paragraph it has
emerged that if we set

\weq{1}{y = t_{n-2}F_0(x) + t_{n-3}F_1(x) \cdots + t_0 F_{n-2}(x)}

then the symmetric functions of the $n$ values $y_1, y_2 \cdots y_n$,
which one obtains from $y$ when one puts for $x$ the $n$ roots of
%CONT>>
